\uuid{YdsH}
\titre{Primitives usuelles}
\auteur{hanani ; mzari}
\organisation{exo7}
\licence{CC BY-NC-SA 4.0}
\datecreate{2020-01-20}
\niveau{L1}
\module{Analyse}
\chapitre{Calcul d'intégrales}
\sousChapitre{Primitives diverses}
\difficulte{1}
\theme{Primitives, primitives usuelles}
\exo7qcmid{lille-2.primitives.003}

\contenu{
\begin{question}
\qclassification{127.11}
\qtags{motcle=primitives usuelles}
Parmi les affirmations suivantes, cocher celles qui sont vraies :
\begin{answers}
\bad{$\displaystyle F(x)=\sqrt{x+1}+\mathrm{e}^{x}$ est une primitive de $\displaystyle f(x)=\frac{1}{\sqrt{x+1}}+\mathrm{e}^{x}$ sur $]-1,+\infty[$.}
\good{$\displaystyle F(x)=2\sqrt{x+1}+\mathrm{e}^{x}$ est une primitive de $\displaystyle f(x)=\frac{1}{\sqrt{x+1}}+\mathrm{e}^{x}$ sur $]-1,+\infty[$.}
\bad{La primitive de $\displaystyle f(x)=\frac{1}{\sqrt{x+1}}+\mathrm{e}^{x}$ sur $]-1,+\infty[$ qui s'annule en $0$ est $\displaystyle F(x)=\sqrt{x+1}+\mathrm{e}^{x}-2$.}
\good{La primitive de $\displaystyle f(x)=\frac{1}{\sqrt{x+1}}+\mathrm{e}^{x}$ sur $]-1,+\infty[$ qui s'annule en $0$ est $\displaystyle F(x)=2\sqrt{x+1}+\mathrm{e}^{x}-3$.}
\end{answers}
\begin{explanations}
Le calcul de $F'(x)$ montre que $\displaystyle 2\sqrt{x+1}+\mathrm{e}^{x}$ est une primitive de $\displaystyle \frac{1}{\sqrt{x+1}}+\mathrm{e}^{x}$ sur $]-1,+\infty[$. Donc, toute primitive de $\displaystyle f(x)$ sur $]-1,+\infty[$ est de la forme : $\displaystyle F(x)=2\sqrt{x+1}+\mathrm{e}^{x}+k$, $k\in \Rr$, et celle qui vérifie $F(0)=0$ est $\displaystyle F(x)=2\sqrt{x+1}+\mathrm{e}^{x}-3$.
\end{explanations}
\end{question}
}
